\section{Related Work}
\label{sec:related}

\subsection{Domain generalization for medical segmentation}
\label{subsec:dg}

Domain generalization (DG) targets unseen distributions without access to
target data, unlike unsupervised domain adaptation, which assumes unlabelled
target images and solves an easier problem~\cite{shin2023sdcuda}. Most
medical DG methods are \emph{multi-source}, interpolating across several
labelled domains by meta-learning, feature alignment, frequency-space
exchange, multi-frequency architectures, foundation-model scale, or
diffusion-based stylization~\cite{li2018mldg,wang2020dofe,liu2021feddg,nam2024madgnet,sengupta2025synthfm,bao2024structaware}
-- all of which still rely on source diversity. We instead work in the
\emph{single-source} regime --- one
labelled dataset, no target data, no second source -- the setting that
matters when target-domain data is unavailable at training time, which is
routine in clinical deployment~\cite{zhang2020bigaug,ouyang2022gin}. This
rules out cross-domain interpolation and makes architecture-coupled methods
orthogonal to our contribution: our mechanism acts at the input level and
leaves the segmentation network a fixed nnU-Net~\cite{isensee2021nnunet}.

\subsection{Synthesis-based contrast-agnostic training}
\label{subsec:synth}

The dominant single-source strategy is label-generative domain randomization,
introduced by SynthSeg~\cite{billot2023synthseg}: intensities are sampled from a
per-label Gaussian mixture over an anatomical segmentation, yielding images of
arbitrary contrast that force contrast-invariance. This family carries two
structural costs: it is \emph{label-driven} (a dense anatomical label map is
required per subject -- SIAM shifts the dependence from label
\emph{quantity} to \emph{quality} with curated templates, but each still
needs a full dense parcellation~\cite{valabregue2026siam}), and it is
\emph{texture-destroying}, since per-label Gaussian sampling replaces real
gradients, partial-volume detail and unlabelled structure with piecewise
noise, retaining layout but not texture. The lineage has been specialized to
white-matter and MS lesions and to atlas-based
expectation-maximization~\cite{laso2024wmhsynthseg,caselles2026timelesseg,cerri2021samseg}.
Among SynthSeg's own variants we take the
original (\emph{noEM}), which resamples intensities from a dense anatomical
label map and never looks at the acquired scan, and the \emph{EM} variant,
which instead fits its per-class intensity model to the real image by
expectation-maximization and so needs only the sparse task labels a
segmentation dataset already provides, not a dense parcellation. The EM
variant has no released implementation; ours follows the mechanism described
in the original paper~\cite{billot2023synthseg}, and a faithful
re-implementation is itself a contribution.

\subsection{Texture as a segmentation signal}
\label{subsec:texture}

That local texture, not only intensity level, carries information a
segmentation model can use is established independently of us: second-order
gray-level statistics discriminate image content beyond simple
intensity~\cite{haralick1973texture}, and quantitative texture features from
tumour regions carry prognostic signal across CT, MRI and
PET~\cite{lambin2012radiomics}. Closer to our setting, Kang et al.\
disentangle structure from texture in cross-domain image translation with a
texture co-occurrence loss, targeting the same failure mode we do --- naive
translation distorts fine texture and degrades downstream
segmentation~\cite{kang2023structurepreserving}; a texture-fidelity loss has
likewise been added to GAN-based lesion augmentation~\cite{guan2022texture}.
That some targets resist boundary-based methods for lack of a sharp edge,
and are better served by texture-based segmentation, has been argued
directly for diffuse white-matter lesions~\cite{kruggel2008texture}. None of
this prior work, to our knowledge, connects texture preservation in an
\emph{augmentation} to \emph{which} segmentation targets it helps; that
connection, tested causally, is our contribution.

\subsection{Single-source intensity and appearance augmentation}
\label{subsec:aug}

Where the first line of single-source methods synthesizes from a label map, a
second manufactures appearance diversity directly from the real image, with
no generative label model. General-purpose libraries stack intensity,
spatial and artifact transforms on an nnU-Net trainer; the recent Auglab
pipeline~\cite{molinier2026auglab,isensee2021nnunet} is one such
GPU-accelerated framework and is the strongest competitor in our
experiments, and BigAug stacks a comparable set of coarse intensity, quality
and spatial perturbations under the same single-source DG protocol we
use~\cite{zhang2020bigaug}. Our method inserts the PALETTE transform as one
additional operator inside the Auglab pipeline, so the comparison against
Auglab directly isolates PALETTE's contribution. Closer to our operator are
random-filter methods, which vary appearance with an untrained nonlinear
filter: GIN applies randomly re-initialized convolutions for label-free
appearance variants~\cite{ouyang2022gin} (the GIN-alone configuration we
compare against is an ablation in the original paper of a two-view
objective on 2-D slices, so our comparison is illustrative of the
random-filter family rather than a reproduction of that paper's proposed
method); related variants establish the random-filter principle for
natural images, perturb intensities per class using the task mask, add
learned diffusion generators, or perturb the Fourier amplitude spectrum
instead of the image
itself~\cite{xu2021randconv,su2023slaug,chen2024diffusionsdg,zhao2024morestyle,li2024raffesdg,qiao2024ras}.
Our operator
instead partitions the foreground into intensity classes and Voronoi
sub-regions and applies a signed \emph{affine} remap of the real intensities
within each, invertible and gradient-preserving by construction. The
closest peer is SRCSM~\cite{thaler2025semrandconv} (semantic-aware random
convolution with source matching), which also localizes its perturbation
region-wise in true 3-D and, like PALETTE, layers exactly one
label-consuming component over an otherwise label-free step: at training
time it draws a fresh untrained convolution per anatomical class from the
ground-truth map (needing a Gaussian smoothing pass to control the
boundary-contrast artifacts this introduces), and at test time applies a
label-free histogram matching onto the source distribution. The difference
from PALETTE is not label use versus none, but \emph{how} that step
transforms the signal: an untrained nonlinear convolution versus a
deterministic affine remap that leaves relative intensity structure inside
each region intact.

\subsection{Cross-contrast evaluation protocol}
\label{subsec:gap}

We evaluate a model trained on a single real input modality on every
held-out modality of the \emph{same} segmentation task, including the
training modality itself, across four tasks spanning brain, abdominal and
tumour anatomy in MRI and CT; a sixth task, spine segmentation, is a
planned extension. This single-sequence-in, every-sequence-out design is
not unprecedented -- Auglab's own evaluation follows the identical protocol
for spine segmentation~\cite{molinier2026auglab} -- and our contribution is
breadth (four distinct cross-contrast tasks, not one) and, for two of them,
a further axis: MS-lesion and abdominal-organ segmentation are each also
evaluated against independent external test datasets acquired at different
sites, so the training dataset's own test set is one point among several
rather than the only held-out evidence. A fifth task, liver-tumour
segmentation, trains on a single modality and so has no cross-contrast axis
of its own, but is evaluated against two independent external cohorts the
same way, extending this same-site-is-not-enough argument to a setting
where cross-contrast transfer cannot be measured at all.
